\documentclass[11pt]{article}
\usepackage[utf8]{inputenc}
\usepackage[margin=1in]{geometry}
\usepackage{xcolor}
\usepackage{enumitem}

\begin{document}

\begin{flushleft}
    {\fontfamily{pzc}\selectfont\Huge\color{magenta}The Laburnum Top}
\end{flushleft}

\begin{flushright}
    {\fontfamily{pzc}\selectfont\LARGE Ted Hughes} \\
    \noindent\rule{1.5in}{0.5pt}
\end{flushright}

\vspace{2em}

\begin{verse}
The Laburnum top is silent, quite still \\
In the afternoon yellow September sunlight, \\
A few leaves yellowing, all its seeds fallen. \\
\end{verse}

\begin{verse}
Till the goldfinch comes, with a twitching chirrup \\
A suddenness, a startlement, at a branch end. \\
Then sleek as a lizard, and alert, and abrupt, \\
She enters the thickness, and a machine starts up \\
Of chitterings, and a tremor of wings, and trillings --- \\
The whole tree trembles and thrills. \\
It is the engine of her family. \\
She stokes it full, then flirts out to a branch-end \\
Showing her barred face identity mask \\
\end{verse}

\begin{verse}
Then with eerie delicate whistle-chirrup whisperings \\
She launches away, towards the infinite \\
\end{verse}

\begin{verse}
And the laburnum subsides to empty.
\end{verse}

\vspace{3em}

\noindent
\colorbox{pink!40}{
\begin{minipage}{\textwidth}
    \vspace{0.5em}
    \begin{description}[leftmargin=6.5em, style=nextline]
        \item[\textbf{laburnum:}] a short tree with hanging branches, yellow flowers and poisonous seeds
        \item[\textbf{goldfinch:}] a small singing bird with yellow feathers on its wings
    \end{description}
    \vspace{0.5em}
\end{minipage}
}

\end{document}
\documentclass[11pt]{article}
\usepackage[utf8]{inputenc}
\usepackage[margin=1in]{geometry}
\usepackage{enumitem}
\usepackage{titlesec}
\usepackage{fancyhdr}

% Header configuration matching page 32 of the textbook
\pagestyle{fancy}
\fancyhf{}
\fancyhead[L]{32}
\fancyhead[R]{\textsc{Hornbill}}
\renewcommand{\headrulewidth}{0pt}

% Configure section headings appearance
\titleformat{\section}{\normalfont\Large\bfseries}{}{0pt}{}
\titlespacing*{\section}{0pt}{1.5em}{1em}

\begin{document}

\section*{Find out}

\begin{enumerate}[label=\arabic*., leftmargin=2em]
    \item What laburnum is called in your language.
    \item Which local bird is like the goldfinch.
\end{enumerate}

\section*{Think it out}

\begin{enumerate}[label=\arabic*., leftmargin=2em]
    \item What do you notice about the beginning and the ending of the poem?
    \item To what is the bird's movement compared? What is the basis for the comparison?
    \item Why is the image of the engine evoked by the poet?
    \item What do you like most about the poem?
    \item What does the phrase ``her barred face identity mask'' mean?
\end{enumerate}

\section*{Note down}

\begin{enumerate}[label=\arabic*., leftmargin=2em]
    \item the sound words
    \item the movement words
    \item the dominant colour in the poem.
\end{enumerate}

\section*{List the following}

\begin{enumerate}[label=\arabic*., leftmargin=2em]
    \item Words which describe `sleek', `alert' and `abrupt'.
    \item Words with the sound `ch' as in `chart' and `tr' as in `trembles' in the poem.
    \item Other sounds that occur frequently in the poem.
\end{enumerate}

\section*{Thinking about language}

\noindent Look for some other poem on a bird or a tree in English or any other language.

\section*{Try this out}

\noindent Write four lines in verse form on any tree that you see around you.

\end{document}
\documentclass[11pt]{article}
\usepackage[utf8]{inputenc}
\usepackage[margin=1in]{geometry}
\usepackage{enumitem}
\usepackage{fancyhdr}
\usepackage{xcolor}

% Header configuration matching page 33 of the textbook
\pagestyle{fancy}
\fancyhf{}
\fancyhead[L]{\textsc{The Laburnum Top}}
\fancyhead[R]{33}
\renewcommand{\headrulewidth}{0pt}

\begin{document}

\begin{center}
    {\color{magenta}\Large\textbf{--- Notes ---}}
\end{center}
\vspace{1em}

\noindent This poem has been placed after a text which has references to names of plants for thematic sequencing.

\vspace{1.5em}
\noindent \textbf{\large Understanding the poem}

\begin{itemize}[leftmargin=1.5em, label=\textbullet]
    \item Glossing of `laburnum' and `goldfinch'
    \item Factual understanding
    \item Movement of thought and structuring (poetic sensitivity)
    \item Focus on figures of speech and imagery used (poetic sensitivity)
    \item Attention to sounds, lexical collocations (poetic sensitivity)
\end{itemize}

\vspace{1.5em}
\noindent \textbf{\large Thinking about language}

\begin{itemize}[leftmargin=1.5em, label=\textbullet]
    \item Finding equivalents in other languages (multilingualism)
    \item Relating to thematically similar poems in other languages (multilingualism)
    \item Attempt at creativity
\end{itemize}

\end{document}

